\section{Rings}
A set $R$ or elements forms a ring if to any two elements $x$ and $y$ of $R$ taken in a particular order there are associated elements of $R$, called their sum and product and denoted by 
$x + y $ and $xy$ respectively, both being unique, and such that axioms  $A1, A2, A3, A4$ and $M2$ and $D1$ above are satisfied. \\

{\bf Basic types of rings}\\
{\twelveit Commutiative rings}\\
The most obvious additional axiom is that of commutativity of multiplication $M!$, that $xy = yx$ for all $x$ and $y$ in the ring. Such a ring is called a {\twelveit commutative ring} (addition is always commutative). \\

{\twelveit Rings with unity}
A ring admits of multiplication but not necessarily of division. For the latter to be possible we must have a neutral element, and also inverses, which are defined in terms of the neutral element. 
If both these are present, the structure becomes a {\twelveit field}. However, a ring may have a neutral element and yet not possess inverses (example -- the ring of integers). Such a rings is called 
a {\twelveit ring with unity}.